In this section we suggest a set of countermeasures that would make the signature scheme more secure against the exposed fault models.

We split the countermeasures into different categories :

\begin{itemize}
    \item Post-publication code path check : same signature size, no modification of the tree construction.
    \item Path construction algorithm change : same signature size.
    \item GGM tree-less signature : larger signature.
\end{itemize}

\subsection{Post-publication path check}

This countermeasure consists is running checks on the output of \SeedTreeToPath.

A simple size check on the path is enough to be secure against fault models 1 and 3, 
since they result in path lengths of $1$ and $\text{golden path size}+1$ respectively.
The expected path size can be computed in linear time from $\calH$. (Todo : algorithm ?).

A more thorough and correct check would be to rebuild the tree using $\PathToSeedTree$ 
on the obtained path and to compare the obtained leaves to $\calH$ and the reference tree.
This can also be done in linear time.

\begin{remark}
    In order to pass this check, it would be sufficient to induce 
    the same fault as in the $\SeedTreeToPath$ on $\PathToSeedTree$, seeing as they both 
    call the same function in the reference implementation.
\end{remark}

\subsubsection{Impact on performances}

We implement both of these checks and compare execution times of the signature step with 
or without countermeasure. The results are summarized in table~\ref{tab:counter_1_cmp}. We refer to the path size check countermeasure by $\textbf{C}_{size}$ and 
the full tree check by $\textbf{C}_{tree}$.

\begin{table}[H]
\centering
\caption{Countermeasures execution time comparison}
\label{tab:counter_1_cmp}
\setlength{\tabcolsep}{10pt}
\begin{tabular}{@{}lcc@{}}
\toprule
\textbf{Version} & \textbf{Time (min)} & \textbf{cycles} \\
 Base   &  0.001230 & 4798061  \\
$C_{size}$   &  0.001070 & 4171382  \\
$C_{tree}$   &  0.001303 & 5080762  \\

\bottomrule
\end{tabular}
\end{table}


\subsection{Path construction algorithm modification}

Though the current implementation of $\SeedTreeToPath\slash\PathToSeedTree$ is very efficient,
it allows for attack namely because it merges the flag tree construction and the path construction
steps. 

Splitting the two steps and introducting a check after both stages could limit the 
attack surface by introducing redundancy.

Todo : detail + implementation + time analysis.

\subsection{GGM-tree-less signature}

We could also consider not using GGM trees and publishing all $t-w$ seeds. 
Todo : size of the signature.